\section{Related Work}
\label{sec:related}

% Compilers then machines. Not a bake-off.

{TVM}, {MLIR}, and {IREE} compile a graph onto a backend that is
already there~\cite{chen2018tvm,lattner2021mlir,iree}.
I needed the other direction: a new {MMIO} engine on a {RISC-V}
{SoC}, named from {PyTorch}, still built by stock {GCC}.

Gemmini generates a systolic array and evaluates it on
FireSim~\cite{genc2021gemmini}.
I start from a finished {SoC} and hang one more engine on the bus,
without a new opcode.

Simba, Occamy, and {SNAX} are the machines: scale-out inference and
high-performance computing across
chiplets~\cite{shao2019simba,bartolini2024occamy,snax}.
This paper is the path from that kind of machine back into a compiler
a {PyTorch} user can run.

FireSim uses {FPGA}s as cycle-approximate {ASIC}
simulators~\cite{karandikar2018firesim}.
I use the {FPGA} as the card: the same {ELF} as the instruction-set
simulator ({ISS}), with boxes back on the host.

{YOLOv3} is the accuracy bar for the running
example~\cite{redmon2018yolov3}.
Matching the integer {mAP} shows the stack closed; it is not a
detector paper.
